% Options for packages loaded elsewhere
\PassOptionsToPackage{unicode}{hyperref}
\PassOptionsToPackage{hyphens}{url}
%
\documentclass[
  11pt,
]{article}
\usepackage{amsmath,amssymb}
\usepackage{iftex}
\ifPDFTeX
  \usepackage[T1]{fontenc}
  \usepackage[utf8]{inputenc}
  \usepackage{textcomp} % provide euro and other symbols
\else % if luatex or xetex
  \usepackage{unicode-math} % this also loads fontspec
  \defaultfontfeatures{Scale=MatchLowercase}
  \defaultfontfeatures[\rmfamily]{Ligatures=TeX,Scale=1}
\fi
\usepackage{lmodern}
\ifPDFTeX\else
  % xetex/luatex font selection
\fi
% Use upquote if available, for straight quotes in verbatim environments
\IfFileExists{upquote.sty}{\usepackage{upquote}}{}
\IfFileExists{microtype.sty}{% use microtype if available
  \usepackage[]{microtype}
  \UseMicrotypeSet[protrusion]{basicmath} % disable protrusion for tt fonts
}{}
\makeatletter
\@ifundefined{KOMAClassName}{% if non-KOMA class
  \IfFileExists{parskip.sty}{%
    \usepackage{parskip}
  }{% else
    \setlength{\parindent}{0pt}
    \setlength{\parskip}{6pt plus 2pt minus 1pt}}
}{% if KOMA class
  \KOMAoptions{parskip=half}}
\makeatother
\usepackage{xcolor}
\usepackage[margin=25mm]{geometry}
\usepackage{color}
\usepackage{fancyvrb}
\newcommand{\VerbBar}{|}
\newcommand{\VERB}{\Verb[commandchars=\\\{\}]}
\DefineVerbatimEnvironment{Highlighting}{Verbatim}{commandchars=\\\{\}}
% Add ',fontsize=\small' for more characters per line
\newenvironment{Shaded}{}{}
\newcommand{\AlertTok}[1]{\textcolor[rgb]{1.00,0.00,0.00}{\textbf{#1}}}
\newcommand{\AnnotationTok}[1]{\textcolor[rgb]{0.38,0.63,0.69}{\textbf{\textit{#1}}}}
\newcommand{\AttributeTok}[1]{\textcolor[rgb]{0.49,0.56,0.16}{#1}}
\newcommand{\BaseNTok}[1]{\textcolor[rgb]{0.25,0.63,0.44}{#1}}
\newcommand{\BuiltInTok}[1]{\textcolor[rgb]{0.00,0.50,0.00}{#1}}
\newcommand{\CharTok}[1]{\textcolor[rgb]{0.25,0.44,0.63}{#1}}
\newcommand{\CommentTok}[1]{\textcolor[rgb]{0.38,0.63,0.69}{\textit{#1}}}
\newcommand{\CommentVarTok}[1]{\textcolor[rgb]{0.38,0.63,0.69}{\textbf{\textit{#1}}}}
\newcommand{\ConstantTok}[1]{\textcolor[rgb]{0.53,0.00,0.00}{#1}}
\newcommand{\ControlFlowTok}[1]{\textcolor[rgb]{0.00,0.44,0.13}{\textbf{#1}}}
\newcommand{\DataTypeTok}[1]{\textcolor[rgb]{0.56,0.13,0.00}{#1}}
\newcommand{\DecValTok}[1]{\textcolor[rgb]{0.25,0.63,0.44}{#1}}
\newcommand{\DocumentationTok}[1]{\textcolor[rgb]{0.73,0.13,0.13}{\textit{#1}}}
\newcommand{\ErrorTok}[1]{\textcolor[rgb]{1.00,0.00,0.00}{\textbf{#1}}}
\newcommand{\ExtensionTok}[1]{#1}
\newcommand{\FloatTok}[1]{\textcolor[rgb]{0.25,0.63,0.44}{#1}}
\newcommand{\FunctionTok}[1]{\textcolor[rgb]{0.02,0.16,0.49}{#1}}
\newcommand{\ImportTok}[1]{\textcolor[rgb]{0.00,0.50,0.00}{\textbf{#1}}}
\newcommand{\InformationTok}[1]{\textcolor[rgb]{0.38,0.63,0.69}{\textbf{\textit{#1}}}}
\newcommand{\KeywordTok}[1]{\textcolor[rgb]{0.00,0.44,0.13}{\textbf{#1}}}
\newcommand{\NormalTok}[1]{#1}
\newcommand{\OperatorTok}[1]{\textcolor[rgb]{0.40,0.40,0.40}{#1}}
\newcommand{\OtherTok}[1]{\textcolor[rgb]{0.00,0.44,0.13}{#1}}
\newcommand{\PreprocessorTok}[1]{\textcolor[rgb]{0.74,0.48,0.00}{#1}}
\newcommand{\RegionMarkerTok}[1]{#1}
\newcommand{\SpecialCharTok}[1]{\textcolor[rgb]{0.25,0.44,0.63}{#1}}
\newcommand{\SpecialStringTok}[1]{\textcolor[rgb]{0.73,0.40,0.53}{#1}}
\newcommand{\StringTok}[1]{\textcolor[rgb]{0.25,0.44,0.63}{#1}}
\newcommand{\VariableTok}[1]{\textcolor[rgb]{0.10,0.09,0.49}{#1}}
\newcommand{\VerbatimStringTok}[1]{\textcolor[rgb]{0.25,0.44,0.63}{#1}}
\newcommand{\WarningTok}[1]{\textcolor[rgb]{0.38,0.63,0.69}{\textbf{\textit{#1}}}}
\setlength{\emergencystretch}{3em} % prevent overfull lines
\providecommand{\tightlist}{%
  \setlength{\itemsep}{0pt}\setlength{\parskip}{0pt}}
\setcounter{secnumdepth}{-\maxdimen} % remove section numbering
\ifLuaTeX
  \usepackage{selnolig}  % disable illegal ligatures
\fi
\usepackage{bookmark}
\IfFileExists{xurl.sty}{\usepackage{xurl}}{} % add URL line breaks if available
\urlstyle{same}
\hypersetup{
  pdftitle={P16.4: Galois descent at the trivial-valuation boundary},
  pdfauthor={Unverified research draft; relative to the supplied complete grant B},
  hidelinks,
  pdfcreator={LaTeX via pandoc}}

\title{P16.4: Galois descent at the trivial-valuation boundary}
\usepackage{etoolbox}
\makeatletter
\providecommand{\subtitle}[1]{% add subtitle to \maketitle
  \apptocmd{\@title}{\par {\large #1 \par}}{}{}
}
\makeatother
\subtitle{A full ample equation over arbitrary characteristic-zero
ground fields, and the unresolved nontrivial-valuation axis}
\author{Unverified research draft; relative to the supplied complete
grant B}
\date{6 October 2026}

\begin{document}
\maketitle

\subsection{Scope and result}\label{scope-and-result}

The entire original P16.4 is not solved in this paper. We prove a
ground-field extension of the supplied complete B theorem at the
\textbf{trivial-valuation boundary}. This result is in every dimension
and on the whole analytic carrier; it is not a projective-line or
skeletal statement. It does not give an operator over a nontrivially
valued nondiscrete field. Accordingly the submission has
\texttt{candidate=false}.

Throughout the positive theorem, \(k\) is any field of characteristic
zero, with its trivial absolute value. Such a field, and its algebraic
closure with the trivial absolute value, are complete. Let \(X\) be a
smooth geometrically integral projective \(k\)-variety of dimension
\(n\geq1\), and let \(L\in\operatorname{Pic}(X)\otimes\mathbb Q\) be an
actual ample rational line bundle. Rational powers are cleared before
using sections. Put \[
 T=X^{\mathrm{an}},\qquad V=(L^n)>0.
\] Geometric integrality is retained to have one generic normalization
and uniqueness modulo one constant. No countability, algebraic
closedness, finite generation of the section ring beyond ampleness,
torus invariance, skeletal support, or restriction on measure support is
imposed. Smoothness, projectivity, algebraizability, ampleness, and
residue characteristic zero are retained. Neither P16.3 nor the nonample
part of P16.5 is asserted.

Use the canonical trivial reference metric on \(L\). Its scalar is
denoted by \(0\). Write \(o\) for the generic trivial valuation. A
basepoint-free Fubini--Study scalar on \(L\) is \[
 h=\frac1m\max_j\bigl(\log\|s_j\|_{\mathrm{triv}}+\lambda_j\bigr),
 \quad s_j\in H^0(X,mL),\quad\lambda_j\in\mathbb R.
 \tag{1}
\] The sections have no common zero. Define \(\operatorname{PSH}_k(L)\)
to consist of their all-point decreasing-net endpoints with finite value
at \(o\). This is the usual fixed ample-bundle psh cone. We prove below
that it equals the fixed-bundle version of the cone in grant B after
scalar extension. Supported points and values \(-\infty\) are included.
Constants mean literal real constants.

For continuous psh \(v\), let
\(\operatorname{MA}_{k,c}(v_1,\ldots,v_n)\) be the genuine continuous
mixed operator, with its raw intersection normalization. Define \[
 \begin{aligned}
 E_k(v)&=\frac1{n+1}\sum_{j=0}^n
 \int_T v\,\operatorname{MA}_{k,c}(v^{[j]},0^{[n-j]}),\\
 E_k(u)&=\inf_{\substack{v\in C(T)\cap\operatorname{PSH}_k(L)\\v\geq u}}E_k(v),
 \qquad\mathcal E^1_k=\{u:E_k(u)>-\infty\}.
 \end{aligned}
 \tag{2}
\] Thus \(E_k(0)=0\) and \(E_k(u+c)=E_k(u)+Vc\). All suprema and
inequalities are on the original full carrier.

\textbf{Theorem 1 (proved boundary theorem).} There are genuine positive
symmetric mixed Radon measures \(\operatorname{MA}_k(u_1,\ldots,u_n)\)
on \((\mathcal E^1_k)^n\), each of raw mass \(V\), extending the actual
continuous mixed operator. They are independent of decreasing continuous
psh approximations, weakly continuous along arbitrary decreasing
finite-energy nets, and satisfy the energy identity, integrability of
every finite-energy test, and its weighted decreasing continuity.

For a positive Radon measure \(\mu\) on \(T\), the exact target
condition is \[
 \mu(T)=V,\qquad
 J_k(\mu):=\sup_{v\in C(T)\cap\operatorname{PSH}_k(L)}
       \left(E_k(v)-\int_T v\,d\mu\right)<\infty.
 \tag{3}
\] It is equivalent to integrating every genuine \(\mathcal E^1_k\)
potential, after any constant normalization. Every such target admits an
actual \(u\in\mathcal E^1_k\) with \(\operatorname{MA}_k(u)=\mu\). Every
solution is unique modulo a real constant. The normalized psh cone and
its bounded energy-cost subsets are compact in divisorial topology.
All-point envelopes, both signs of obstacle energy differentiation,
strict Borel locality, comparison, and actual variational attainment
hold. No inverse construction is used to define the operator.

The proof is relative only to the \textbf{complete B body supplied in
the assignment}, in its algebraically closed trivially valued smooth
projective actual-big rational-bundle domain, specialized to ample
bundles. Its internal G0/G1 references are not separately imported.
Historical review labels, old notes and old source hashes are not
premises. In particular, the unreviewed round46 projective-line operator
is not assumed. Ordinary field algebra, intersection theory,
compact-space Radon measures, and the continuous-metric theory are used
where specified. All new descent statements are proved here.

\subsection{Exact use of B and the continuous
background}\label{exact-use-of-b-and-the-continuous-background}

Let \(\bar k\) be an algebraic closure, \(Y=X_{\bar k}^{\mathrm{an}}\),
\(G=\operatorname{Gal}(\bar k/k)\), and \(\pi:Y\to T\). The full B
theorem applies on \(X_{\bar k}\) with the actual ample \(L_{\bar k}\).
Its reference is \(R=0\): zero is the scalar of a generated power of
\(L_{\bar k}\), and every normalized section scalar has supremum zero at
the generic trivial valuation. Thus the B target condition has no
reference-pole exception and is \[
 D_{\bar\nu}=\sup_{u\in\mathcal E^1_{\bar k},\,u(o)=0}
       \left(E_{\bar k}(u)-\int_Yu\,d\bar\nu\right)<\infty.
 \tag{4}
\] B supplies actual psh closure and envelopes, divisorial compactness
and all-point determination, scalar energy and its decreasing
continuity, genuine mixed products and weighted continuity, obstacle
differentiation, comparison, intrinsic targets, attainment, and
uniqueness for every solution. No version of these assertions over \(k\)
or a nontrivially valued field is granted.

For the genuine continuous background we checked the fixed primary
version Boucksom--Eriksson, arXiv:1805.01016v3: Definition 7.2,
Corollary 7.6, Theorem 7.14, Proposition 8.3 including its proof,
Theorem 8.4 with Lemma 8.6 and the full extension proof, and Definition
9.13--Proposition 9.14 including the proof. These units allow complete
fields, including trivial valuation, and continuous psh metrics on ample
rational bundles. Their field-extension law is a \textbf{pushforward}
law with unchanged geometric raw mass, not multiplication by the degree
of the field extension. They supply no singular field-extension law;
that law is proved below. The compact Hausdorff and seminorm
descriptions of projective analytifications are standard analytic
background, also recorded in its Section 4.1.

Here is a useful bridge between B's varying-bundle cone and (1). Suppose
\(h_i\downarrow u\) is a defining net in B, on actual \(B_i\) with
\(c_1(B_i)\to c_1(L)\), and \(u(o)\) finite. On a tail take
\(h_i\leq C\). For rational \(0<t<1\) and late \(i\), \[
 L-(1-t)B_i\text{ is ample}.
\] Tensor the section system for \((1-t)h_i+tC\) with a weight-zero
generated power of this difference; the scalar is unchanged everywhere.
Indeed at the valuation-ring center of each analytic point some
generator is a unit. Such centers exist by projectivity, including for
supported semivaluations. Append a generated power of \(L\) at scalar
\(C-N\). This gives the basepoint-free scalar \[
 H_{i,t,N}=\max\bigl((1-t)(h_i-C)+C,C-N\bigr)
 \quad\text{on the actual fixed }L.
 \tag{5}
\] Eligible triples are directed by later \(i\), smaller \(t\), and
larger \(N\). A common later \(i\) fits the prescribed positive ample
differences. Since \(h_i-C\leq0\), the displayed functions decrease in
that order. Their infimum is \(u\) at every point: first hold \(t,N\)
and take \(i\), then let \(t\downarrow0\) and \(N\uparrow\infty\). At a
pole this is a literal limit of finite upper bounds, with no subtraction
of infinities. Conversely a fixed-\(L\) net is already a defining net in
B. This proves equality of the cones in the ample subcase.

B's genuine operator agrees there with the continuous operator just
described. Rational weighted FS systems have actual finite intersection
measures; these are the same finite coefficients in B. Real weights are
approximated uniformly by rational weights. A continuous psh metric is
uniformly approximated by such systems; if the error of the \(j\)th is
at most \(4^{-j}\), adding \(2\cdot4^{-j}\) makes them decrease to the
metric. B's decreasing continuity and the continuous operator's uniform
continuity identify both limits. This establishes agreement without
postulating a new unbounded product.

\subsection{The analytic quotient, including all
fibers}\label{the-analytic-quotient-including-all-fibers}

\textbf{Lemma 2 (quotient).} The action of the profinite group \(G\) on
\(Y\) is jointly continuous; \(\pi\) is surjective and its fibers are
exactly the \(G\)-orbits. It induces a homeomorphism \(Y/G\simeq T\).
The point over \(o\) is the generic trivial valuation of \(X_{\bar k}\).

\textbf{Proof.} On an affine \(k\)-chart, an analytic point is a
multiplicative seminorm on its coordinate algebra. Every function after
scalar extension has coefficients in a finite extension. Its conjugates
therefore vary locally constantly with \(g\in G\), proving joint
continuity of the action in the pointwise seminorm topology.

For a finite Galois extension \(E/k\) and \(x\in T\), its analytic fiber
is the set of multiplicative norms on the finite étale algebra
\(\mathcal H(x)\otimes_k E\). Here \(\mathcal H(x)\) is the completed
residue field; it is complete and hence henselian. The algebra is a
nonzero product of finite separable fields, each with the unique
extension of the absolute value of \(\mathcal H(x)\). This also holds
when that absolute value is trivial. Thus the fiber is nonempty.
\(\operatorname{Gal}(E/k)\) acts transitively on its factors: otherwise
the sum of the idempotents in one orbit would be a nontrivial invariant
idempotent, whereas the invariant algebra is \(\mathcal H(x)\).

Analytifications over \(\bar k\) are the inverse limit of the
finite-extension analytifications. To check this directly, choose a
\(k\)-affine chart containing the scheme support of \(x\). A compatible
family of seminorms defines a seminorm on the union of its
scalar-extended algebras; all algebraic norm identities are
finite-coefficient identities. This is the required point over
\(\bar k\), and the seminorm topologies give the claimed inverse-limit
topology. Surjectivity follows from the nonempty finite fibers and
compactness of their inverse system. For two points \(y,z\) in the same
fiber, the sets \[
 \{g\in G:g y_E=z_E\}
\] are nonempty closed subsets of \(G\). Any finite intersection is
nonempty by passing to a finite Galois compositum. Compactness of \(G\)
gives one \(g\) matching every finite level. Consequently \(gy=z\).

Both \(Y\) and \(T\) are compact Hausdorff. The continuous bijection
\(Y/G\to T\) is therefore a homeomorphism. At the generic trivial
valuation, geometric integrality makes \(k(X)\otimes_k E\) a field for
each finite \(E\). Its only extension of the trivial absolute value is
trivial, so this fiber is a singleton. No points with nongeneric support
were omitted. \(\square\)

An invariant usc extended-real function on \(Y\) consequently descends
to an usc function on \(T\). For example, its closed upper level sets
have compact, hence closed, images. An invariant continuous function
descends continuously. This remains valid at \(-\infty\).

\subsection{Exact descent of actual FS systems and singular
potentials}\label{exact-descent-of-actual-fs-systems-and-singular-potentials}

\textbf{Lemma 3 (FS descent).} For any FS scalar \(h\) on
\(L_{\bar k}\), its orbit maximum \[
 h^G(y)=\max_{g\in G}h(gy)
 \tag{6}
\] is the pullback of an actual FS scalar on \(L\) with the same cleared
degree.

\textbf{Proof.} All sections in \(h\) are defined over one finite Galois
extension \(E/k\). In degree \(m\), take the finite union of all
conjugate section lists, with unchanged weights. This produces \(h^G\).
For each distinct weight \(a\), the subspace \(W_a\) spanned over \(E\)
by sections of weight at least \(a\) is semilinearly Galois stable. It
descends to a \(k\)-subspace \(U_a\): in a fixed \(k\)-basis of
\(H^0(X,mL)\), the unique reduced-row-echelon matrix of \(W_a\) is fixed
by every automorphism, so its entries lie in \(k\). This also proves
\(W_a=U_a\otimes_kE\).

Choose a \(k\)-basis adapted to the finite nested flag \(U_a\), and give
each basis vector its largest flag weight. The weighted maximum of this
basis equals the weighted maximum of the original conjugate generators.
For one inequality express an adapted vector in the generators of its
weight level; for the other express each generator in adapted vectors of
at least its weight. Every nonzero coefficient has absolute value one.
The ultrametric inequality proves both bounds on the section norms at
every analytic point, even when a section vanishes. Basepoint freeness
descends faithfully flatly. Hence this is an actual FS system over
\(k\), and not just a function on a quotient. \(\square\)

\textbf{Lemma 4 (psh descent).} Pullback gives a bijection \[
 \operatorname{PSH}_k(L)\ \simeq\
 \operatorname{PSH}_{\bar k}(L_{\bar k})^G
 \tag{7}
\] as actual functions on all points. It also identifies continuous psh
metrics.

\textbf{Proof.} Pulling back sections proves one direction. If an
invariant \(u\) is a B-psh endpoint, (5) supplies a decreasing net of
fixed-bundle basepoint-free FS scalars \(h_i\downarrow u\). Their orbit
maxima \(h_i^G\) decrease, and Lemma 3 descends each one. Their endpoint
is exactly \(u\). Indeed fix \(y\) and a finite real threshold
\(a>u(y)\). The compact sets \[
 \{g\in G:h_i(gy)\geq a\}
\] decrease to the empty set, since \(u(gy)=u(y)\). Finitely many
suffice to have empty intersection, and directedness supplies a whole
tail with empty set. This proves \(\inf_i h_i^G(y)\leq u(y)\); the
reverse inequality is immediate. If \(u(y)=-\infty\), apply this
argument at each finite threshold. There is no measurable-net
integration or interchange of a noncompact supremum and infimum.

The descended functions decrease all-point to the descended \(u\), whose
value at \(o\) is finite. If \(u\) is continuous, net Dini follows from
the same nested compact-superlevel argument on \(Y\); hence the
descended FS net converges uniformly. Conversely continuous pullback is
continuous. \(\square\)

This proof depends on the trivial coefficient absolute value and on
having an actual upstairs psh potential. It does not provide an upstairs
finite-energy operator for a nontrivially valued field.

\subsection{Radon measures: the unique invariant
lift}\label{radon-measures-the-unique-invariant-lift}

\textbf{Lemma 5.} Every positive Radon measure \(\mu\) on \(T\) has a
unique \(G\)-invariant positive Radon lift \(\widetilde\mu\) on \(Y\).
It satisfies \[
 \pi_*\widetilde\mu=\mu,\qquad \widetilde\mu(Y)=\mu(T),
 \quad\int_Y(f\circ\pi)d\widetilde\mu=\int_T f\,d\mu
 \tag{8}
\] for every nonnegative Borel \(f\), and for every integrable signed
\(f\). The lift and pushforward are inverse homeomorphisms between base
measures and invariant measures for the weak topologies.

\textbf{Proof.} The probability Haar functional on \(G\) can be
constructed by uniform averaging on finite quotient groups; locally
constant functions are uniformly dense on the profinite compact group.
For \(f\in C(Y)\) put \[
 Af(y)=\int_G f(gy)\,dg.
\] Joint continuity and compactness make \(Af\) continuous and
invariant. It therefore descends to a continuous function on \(T\). The
positive functional \(f\mapsto\int_T Af\,d\mu\) has norm \(\mu(T)\);
Riesz representation defines \(\widetilde\mu\). Averaging leaves a
pullback \(f\circ\pi\) unchanged, proving its pushforward and mass.
Group invariance follows from Haar invariance. If \(\lambda\) is another
invariant Radon measure, finite-average approximation, uniform in \(y\),
gives \(\lambda(f)=\lambda(Af)\). Its pushforward determines the latter,
so \(\lambda=\widetilde\mu\). The same continuous tests prove both weak
continuity statements. The Borel integral laws now follow from the usual
pushforward law for a single measure, with countable simple-function
approximation. No arbitrary measurable-net convergence is used.
\(\square\)

On a finite fiber with \(r\) points, its orbit probability puts mass
\(1/r\) at each point. In general \(r\) need not equal \([E:k]\). A
uniform division by the field degree at every point is wrong. We
preserve the original measure, with no mass discarded at supported
points.

\subsection{Raw energy and genuine mixed-operator
descent}\label{raw-energy-and-genuine-mixed-operator-descent}

The finite genuine model normalization deserves emphasis. For a rational
weighted FS test configuration over \(k[t]\), a central prime \(Q\) of
multiplicity \(b_Q\) contributes \[
 b_Q\bigl(\mathcal L_1|_Q\cdots\mathcal L_n|_Q\bigr)_k\,
       \delta_{b_Q^{-1}\operatorname{ord}_Q|_{k(X)}}.
 \tag{9}
\] Constant finite separable extension splits primes but does not ramify
their local uniformizers. A finite étale extension of a DVR is
unramified. Intersection degrees over \(E\) obey \[
 \deg_k(p_*z)=[E:k]\deg_E(z).
\] The projection formula for the finite map on the total model
consequently cancels this same factor \([E:k]\). Summing over split
components gives exactly (9). In particular \((L_E^n)_E=(L^n)_k=V\);
this is geometric base change, not a degree-\([E:k]\) map measured
everywhere over \(k\). Real weights pass by uniform approximation. For
continuous psh metrics the same pushforward formula follows from uniform
FS approximation and the continuous operator. Equivalently it is
Proposition 8.3(iv) passed through Theorem 8.4 in the checked version.
Therefore \[
 E_k(v)=E_{\bar k}(v\circ\pi)
 \quad(v\text{ continuous psh}).
 \tag{10}
\]

Both the continuous energy and operator over \(\bar k\) are invariant
under the semilinear \(G\)-action. At rational FS stages this follows by
transporting the actual section system and finite intersections; all
degrees and multiplicities are preserved. Uniform limits prove it on
continuous metrics. Decreasing approximants and B prove it on the entire
finite-energy space and its mixed products.

\textbf{Lemma 6 (extended energy).} For every actual
\(u\in\operatorname{PSH}_k(L)\), \[
 E_k(u)=E_{\bar k}(u\circ\pi).
 \tag{11}
\] This includes the value \(-\infty\). Thus (7) identifies
\(\mathcal E^1_k\) with the invariant part of \(\mathcal E^1_{\bar k}\).

\textbf{Proof.} Every continuous base majorant pulls back to a
continuous upstairs majorant, giving the inequality
\(E_k(u)\geq E_{\bar k}(u\circ\pi)\). Lemma 4 supplies actual descended
continuous FS majorants \(a_i\downarrow u\). Their energies equal the
upstairs energies by (10); B's decreasing energy continuity makes these
tend to \(E_{\bar k}(u\circ\pi)\). They are competitors in (2), proving
the reverse inequality. This is cofinality in energy, not an assertion
that averaging a majorant lowers its energy. \(\square\)

Define, before considering any target, \[
 \operatorname{MA}_k(u_1,\ldots,u_n)
 :=\pi_*\operatorname{MA}_{\bar k}(u_1\circ\pi,\ldots,u_n\circ\pi).
 \tag{12}
\] This uses B's genuine operator in its granted domain, on actual
pullback psh metrics. Its continuous restriction equals the genuine
continuous operator by (9)--(10), so it is an actual extension and not
an assigned inverse. Positivity, symmetry and raw mass \(V\) follow
immediately. For any independent decreasing continuous psh approximants
with finite-energy endpoints, their pullbacks are genuine approximants
upstairs, and B's independence and decreasing continuity prove
independence and continuity of (12) along the original directed nets.

The upstairs mixed measure in (12) is invariant and hence is precisely
the invariant lift of the base mixed measure, by Lemma 5. All genuine
integrability and weighted continuity therefore descend by the exact
pushforward integral law. In particular \[
 E_k(u)-E_k(v)=\frac1{n+1}\sum_{j=0}^n
    \int_T(u-v)\,d\operatorname{MA}_k(u^{[j]},v^{[n-j]}).
 \tag{13}
\] The signed difference is taken on the common full-measure set where
both finite-energy tests are finite. No infinite arithmetic at supported
poles is used. Constant shifts preserve the operator and add \(Vc\) to
the energy.

The space and measure condition are independent of the chosen continuous
semipositive reference. If the new reference has scalar \(a\), the
continuous energy cocycle and the bijection of absolute continuous
majorants give \(E_a(u-a)=E_0(u)-E_0(a)\) on the entire cone. Therefore
\(J_a(\mu)=J_0(\mu)-E_0(a)+\int a\,d\mu\); its correction is finite
because \(a\) is continuous and the mass is \(V\). Absolute metrics and
their operator are unchanged.

These definitions are independent of the chosen algebraic closure: an
isomorphism between two algebraic closures transports actual section
systems and invariant measures, and the continuous and decreasing
constructions commute with it. For any algebraic extension \(F/k\)
within \(\bar k\), the same proof with
\(H=\operatorname{Gal}(\bar k/F)\) gives \[
 E_F(u_F)=E_k(u),\quad
 p_*\operatorname{MA}_F(u_{1,F},\ldots,u_{n,F})
       =\operatorname{MA}_k(u_1,\ldots,u_n).
 \tag{14}
\] The measure upstairs is obtained by pushing the \(G\)-invariant lift
to \(X_F^{\mathrm{an}}\); for finite Galois \(F/k\) it is the unique
invariant lift under that finite group. For non-Galois \(F\), this
specifies the lift without pretending that its fibers are orbits of
\(\operatorname{Aut}(F/k)\). Characteristic zero makes every algebraic
extension separable; no inseparable descent is tacitly asserted.

\subsection{Compactness, envelopes, and differentiation on the
base}\label{compactness-envelopes-and-differentiation-on-the-base}

\textbf{Divisorial topology.} Every rationally scaled divisorial
valuation upstairs is defined on a model over some finite \(E/k\) and
restricts to a rationally scaled divisorial valuation of \(k(X)\). One
way to see this is to use the finite extension \(E(X)/k(X)\):
restriction preserves residue transcendence degree \(n-1\) and a
discrete rank-one value group. The associated divisorial valuation has a
proper algebraic model, and characteristic-zero resolution realizes it
on a smooth model over \(k\). Conversely a divisorial valuation over
\(k\) extends divisorially after constant extension: its DVR tensored
with \(E\) is finite étale, and each local prolongation is a DVR with
ramification index one. A finite model defining an upstairs prime
suffices, even over an uncountable \(\bar k\). These facts identify the
coordinate topologies in (7).

B's normalized compact psh cone has closed invariant locus. Invariance
is the collection of closed coordinate equalities \(u(z)=u(gz)\) at
upstairs divisorial coordinates; B's all-point determination promotes
those equalities to invariance at every supported point. Generic
normalization is preserved. Thus the original normalized cone is
compact, and energy weak upper semicontinuity from (11) makes each
normalized cost set \(-E_k(u)\leq C\) compact. Genuine decreasing
closure follows by applying B upstairs and then Lemma 4. This is
potential compactness, not just compactness of measures.

\textbf{Envelopes.} Let \(f\) be an upper bounded usc obstacle on \(T\)
with a genuine psh minorant. Its pullback is invariant. The upstairs raw
all-point envelope is invariant, since \(G\) permutes its competitors.
Lemma 4 descends it. Maximality in each direction gives \[
 P_k(f)\circ\pi=P_{\bar k}(f\circ\pi).
 \tag{15}
\] Both sides are genuine psh functions at all points. Upper-envelope
closure and its usc regularization descend too; quotient maps by compact
groups are open (the saturation of an open set is open), which also
checks compatibility of regularization. At continuous obstacles the
ample B theory gives continuity: in B one may choose the Kodaira
decomposition \(L=1\cdot L+0\), so its continuity locus \(U\) is the
entire carrier for model obstacles; its envelope sup-norm Lipschitz law
and uniform model density give the continuous-obstacle case. Thus no
non-Noetherian regular model is being hypothesized here. Every actual
finite construction is over the Noetherian characteristic-zero ring
\(k[t]\), with ordinary equivariant algebraic resolution; its extension
to \(\bar k\) remains in precisely B's field domain.

\textbf{Both differentiation signs.} For \(u\in\mathcal E^1_k\) and
\(h\in C(T)\), (11), (15), and B's outer derivative give \[
 \left.\frac d{dt}\right|_{t=0}E_k(P_k(u+th))
       =\int_T h\,d\operatorname{MA}_k(u).
 \tag{16}
\] The envelope lies between \(u\pm |t|\|h\|\), so it is actually finite
energy. Both signs are inherited from the proved B derivative, whose
application has the same smooth projective trivially valued
characteristic-zero ample domain. Affine finite-energy differentiation
can also be checked directly: set \(u_t=(1-t)u+tv\) in (13), divide by
\(t\), and expand the mixed slots. B's genuine mixed polynomial law,
preserved by (12), makes every fixed-test integral a polynomial in
\(t\), with finite coefficients by integrability. At \(t=0\) all \(n+1\)
terms are the same integral, so the right derivative is
\(\int(v-u)\,d\operatorname{MA}_k(u)\); the other endpoint and interior
parameters follow identically. This uses the actual mixed map, not just
continuous volume differentiation.

The actual orthogonality assertion follows as well. For \(f\in C(T)\)
put \(a=P_k(f)\), which is continuous by the envelope argument. For
\(0\leq t\leq1\), the obstacle \(a+t(f-a)\) lies between \(a\) and
\(f\), so its psh envelope is exactly \(a\) by maximality. Apply the
right derivative in (16), with the fixed continuous \(h=f-a\), to this
constant-energy path. It gives \[
 \int_T(f-P_k(f))\,d\operatorname{MA}_k(P_k(f))=0.
\] The test is an actual original continuous function and the measure is
the genuine operator already constructed.

\textbf{Locality and comparison.} B's strict Borel locality applied to
the pullbacks pushes to the identical winner sets on \(T\). Hence \[
 \mathbf1_{\{u>v\}}\operatorname{MA}_k(\max(u,v))
       =\mathbf1_{\{u>v\}}\operatorname{MA}_k(u),
\] and the symmetric identity holds. Its comparison principle gives \[
 \operatorname{MA}_k(v)\{u<v\}\leq
       \operatorname{MA}_k(u)\{u<v\}.
 \tag{17}
\] Winner sets are Borel; equality of pullback sets is literal at
\(-\infty\) as well. The null exceptional sets for integrable
finite-energy tests are inherited without restriction or renormalization
of any measure.

\subsection{Intrinsic targets, actual attainment, and every-solution
uniqueness}\label{intrinsic-targets-actual-attainment-and-every-solution-uniqueness}

The continuous-test definition (3) equals the supremum over all
finite-energy tests. Indeed choose any all-point descending continuous
psh majorants \(v_i\downarrow u\in\mathcal E^1_k\). On a tail
\(v_i\leq C\). The nonnegative continuous functions \(C-v_i\) increase
to \(C-u\), and their directed integrals against a Radon measure
converge to the integral of this supremum. Here is the required directed
argument: for each strict superlevel of the supremum, its
continuous-stage superlevels cover it; an inner compact subset admits a
finite subcover, and one directed index dominates that cover. Inner
regularity, finite level partitions and truncation give the integral
equality. This proves the assertion without measurable-net monotone
convergence. Decreasing energy continuity then passes
\(E_k(v_i)-\int v_i\,d\mu\) to the corresponding extended objective. If
\(J_k\) is finite, the integral of every normalized \(-u\) is finite.
Continuous tests are already finite-energy tests, so the two suprema are
equal.

The same reasoning in the ample B domain identifies (4) with its
continuous-test supremum. The decisive equality is \[
 J_k(\mu)=J_{\bar k}(\widetilde\mu).
 \tag{18}
\] Pullback of base continuous tests and (10) give the inequality
\(\leq\). For the reverse take any continuous psh \(v\) upstairs and
average the \textbf{scalar} \(v\) over \(G\): \[
 Av=\int_G v\circ g\,dg.
\] This is a continuous psh function: Haar integrals are uniformly
approximated by finite convex combinations of conjugates, all on the
same actual \(L_{\bar k}\), and the continuous psh cone is closed under
such uniform limits. Finite concavity, Galois invariance of energy, and
its sup-norm continuity give \(E_{\bar k}(Av)\geq E_{\bar k}(v)\).
Invariance of \(\widetilde\mu\) gives
\(\int Av\,d\widetilde\mu=\int v\,d\widetilde\mu\). By Lemma 4, \(Av\)
descends to a continuous base psh test. Its base objective is at least
the original upstairs objective, proving (18). Only continuous functions
are averaged; no averaging or subtraction of unbounded poles occurs.

Consequently \(J_k(\mu)<\infty\) is exactly B's intrinsic target
condition for \(\widetilde\mu\). To check the converse intrinsic
characterization using only original base tests, suppose every
normalized base E1 potential is \(\mu\)-integrable. Write
\(B_k(u)=-E_k(u)\) and \(I_\mu(u)=\int(-u)\,d\mu\). Their integrals are
uniformly bounded on \(B_k(u)\leq1\). Otherwise choose normalized
\(u_j\) there with \(I_\mu(u_j)\geq4^j\) and form \[
 p_m=\sum_{j=1}^m2^{-j}u_j+2^{-m}0.
\] This is a genuine decreasing base psh sequence with generic value
zero and cost at most one by concavity. Base decreasing closure and
energy continuity give \(p\in\mathcal E^1_k\). Ordinary sequential
monotone convergence makes \(I_\mu(p)=\infty\), a contradiction. Let
\(C\) be the uniform bound. The actual contraction proved in B,
transported by (11), is \[
 B_k(ru)\leq C_0r^2B_k(u),\qquad C_0=2^{n+1},\quad0\leq r\leq1.
\] For \(B_k(u)>1\) take \(r=(C_0B_k(u))^{-1/2}\); then \(ru\) has cost
at most one, so \(I_\mu(u)\leq C\sqrt{C_0B_k(u)}\). For cost at most one
use \(C\) directly. Therefore \(\sup_u(I_\mu(u)-B_k(u))\) is finite,
proving \(J_k(\mu)<\infty\). This argument averages no singular
functions and uses only a countable convex series. No support
restriction is imposed by the lifting operation.

Apply B to this target. It supplies an actual upstairs solution
\(\bar u\in\mathcal E^1_{\bar k}\), normalized at the fixed generic
trivial point. For every \(g\in G\), the transformed \(\bar u\) solves
the same invariant target equation. B's uniqueness for every solution
gives equality modulo a constant; generic normalization makes that
constant zero. Thus \(\bar u\) is invariant. Lemmas 4 and 6 descend it
to an actual original \(u\in\mathcal E^1_k\). Formula (12) gives
\(\operatorname{MA}_k(u)=\mu\).

This also proves actual variational attainment. The normalized solution
upstairs is a maximizer by B's whole supporting plane; its descended
objective is identical by (8) and (11). Every base competitor pulls back
to a permitted upstairs competitor. Alternatively B's intrinsic-target
integral continuity on compact normalized energy sublevels restricts
through (7), and its coercivity and weak energy upper semicontinuity
give a base maximizing subnet with an actual E1 endpoint. At such an
endpoint, (16) and both signs infer the equation against every
continuous original test. The operator had already been constructed in
(12).

Finally let \(u,v\) be any two base solutions, without assuming either
was produced by this construction. Their upstairs operator measures are
invariant, and their pushforwards both equal \(\mu\). Lemma 5 forces
both upstairs measures to equal \(\widetilde\mu\). Hence their pullback
potentials are two genuine B solutions and differ by one real constant.
Surjectivity of \(\pi\) proves \(u-v\) is that same constant on every
original point. This is every-solution uniqueness.

Conversely every output of (12) satisfies (3), since its upstairs output
is a B intrinsic target and (18) applies. This completes Theorem 1 in
its exact boundary domain. In dimension zero geometric integrality gives
\(X=\operatorname{Spec}k\): the mass is \(1\), potentials are constants,
and the point measure is the unique target; the same conclusion is
immediate.

\subsection{A concrete non-toric, non-algebraically-closed
application}\label{a-concrete-non-toric-non-algebraically-closed-application}

Give \(k=\mathbb Q\) the trivial absolute value. Let \(C\) be the smooth
elliptic curve \(y^2z=x^3-xz^2\), with origin \(O=[0:1:0]\), and set
\(X=C\times C\). It is a smooth geometrically integral projective
surface. Its nonzero regular one-forms exclude smooth projective toric
surfaces, which are rational; thus it is not toric. Take the actual
ample line bundle
\(L=p_1^*\mathcal O_C(3O)\otimes p_2^*\mathcal O_C(3O)\). If \(F_1,F_2\)
are the two fiber classes, then \(F_1^2=F_2^2=0\) and \(F_1F_2=1\), so
\(V=L^2=18\).

Let \(z\) be the divisorial valuation of the exceptional prime of the
blowup of \(X\) at \((O,O)\). The target \(18\delta_z\) has finite
intrinsic dual energy. Indeed normalized potentials have \(u\leq0\),
\(E_k(u)\leq0\), and \(u(z)\geq-C_z\) for one finite coordinate bound
\(C_z\): extend \(z\) divisorially to \(\bar k\) and use B's compact
ample coordinate bound, then descend. Hence every continuous-test
objective is at most \(18C_z\). The theorem therefore gives an actual
finite-energy metric with raw Monge--Ampère measure \(18\delta_z\), and
every solution differs from it by a constant. This is an application on
an original non-algebraically-closed ground field in dimension two. It
remains a \textbf{trivially valued} application, not a nondiscrete
nontrivial-valuation theorem.

\subsection{What this replaces, and what remains
unproved}\label{what-this-replaces-and-what-remains-unproved}

For the \textbf{proved trivial-valuation boundary}, normalized potential
compactness is replaced by the closed invariant locus in B's actual
compact cone, with divisorial restrictions checked above. Envelope
regularity is replaced by the actual all-point equality (15) and FS
descent. Uniform approximation and energy orthogonality use B only over
\(\bar k\) and pass through (10)--(16). Genuine mixed and unbounded
products are (12), identified with continuous intersections before any
equation is considered; weighted limits use exact pushforward
integration. Scalar and obstacle derivatives are (13) and (16). Borel
locality and comparison are (17). Target integrability, attainment, and
every-solution uniqueness are proved in the preceding section. Every
field extension used is algebraic and trivially valued. No
non-Noetherian valuation ring occurs.

For the original requested \textbf{nontrivially valued arbitrary
complete residue-characteristic-zero field} and arbitrary-dimensional
smooth projective geometrically integral \(X\) with actual ample
rational \(L\), the following statements have not been proved in this
submission:

\begin{enumerate}
\def\labelenumi{\arabic{enumi}.}
\tightlist
\item
  Actual normalized psh compactness and all-point divisorial or
  equivalent determination over nondiscrete fields, strong enough to
  produce original potentials rather than merely limits of measures.
\item
  Actual envelope and regularization inputs over their valuation rings,
  including the uniform approximation or orthogonality link needed for
  unbounded potentials. Characteristic-zero resolution of algebraic
  \(k[t]\) models is not resolution over an arbitrary non-Noetherian
  valuation ring.
\item
  Approximation-independent genuine mixed/unbounded products on the
  entire E1 space, positivity, no loss of the full raw mass \(V=(L^n)\),
  and all required unbounded test integrability and weighted limits.
\item
  Scalar and obstacle energy differentiation against all continuous
  tests for that actual operator, with both signs of any perturbation
  and without replacing it by a completed measure inverse.
\item
  Comparison or equivalent finite-energy estimates sufficient for
  every-solution uniqueness, and actual variational attainment in the
  original potential space for every mass-\(V\) target of finite
  intrinsic dual energy.
\item
  If completed algebraic closure is used, all singular metric, energy
  and measure lifting/descent laws over that nontrivial valuation. Its
  closure is still nontrivially valued and is outside grant B. A
  profinite action on an algebraic closure does not itself supply its
  missing finite-energy theory, and a completed extension requires its
  own topological and geometric checks.
\end{enumerate}

These are specific gaps of this submission, not impossibility theorems
or a claim of exhaustive knowledge of post-2020 literature. The intended
target condition in that unresolved domain remains the raw mass \(V\)
and finite supremum
\(\sup_{v\,\mathrm{continuous\ psh}}(E(v)-\int v\,d\mu)\) relative to a
genuine continuous semipositive reference. Merely defining this supremum
or E1 does not construct the operator. Our theorem does not solve the
nontrivial-valuation part of P16.4 and does not earn
\texttt{candidate=true}.

\subsection{Source audit, provenance, and reproducible
evidence}\label{source-audit-provenance-and-reproducible-evidence}

The precise application ledger is
\texttt{source-application-ledger.json} with a readable companion
\texttt{source-application-ledger.md}. It distinguishes the complete
supplied B grant, independently read continuous-theory proof units,
diagnostic scope checks, and ordinary background reconstructed here. The
old projective-line paper and its three note fields have no registered
support IDs and are not used as theorem grants. We make no claim to have
repaired every part of that old tree theorem; the present positive
theorem does not depend on it.

Fresh retrieval of arXiv:1204.6277v2 reports \textbf{608 PDF pages},
agreeing with the assignment's actual-page statement, rather than the
old metadata figure 591. Lemma 8.5.5 and Proposition 8.7.2, including
their complete pertinent proofs, retain continuous multipliers and local
uniform convergence. This documentary check supplies no unbounded
equation theorem. Pille--Schneider arXiv:2209.04879v2 Section 4.1,
Theorem 4.3 and Remark 4.4 distinguish the continuous extension from the
general finite-energy operator issue. The synthetic paper
arXiv:2307.01697v1 expressly separates its measure results from
investigating actual finite-energy potentials. None is promoted into a
missing nondiscrete operator.

The preserved retrieval objects and rendered text have SHA-256 records.
Those hashes bind \textbf{tool-response snapshots}, not downloaded PDF
bytes: direct \texttt{curl} retrieval failed at the configured proxy,
and no proxy or runtime setting was changed. One ref-ID batch resolved
to a different document; the ledger marks that batch diagnostic-only and
uses subsequently retrieved literal version URLs for the needed units.
The arXiv HTML renderer has a later rendering date than the fixed
version's submission date; the fixed-version PDF was cross-checked at
the applied continuous-law units. No source ledger is treated as a
proof.

\texttt{finite\_descent\_checks.py} checks the new mass and
symmetrization bookkeeping with rational arithmetic on one specified
five-node tree and a reflection group. It enumerates finite grids,
compares kernel and edge energies, checks orbit/stabilizer masses,
Jensen improvement of the objective at invariant targets, and the exact
unbounded-ray truncation formulas. These are finite model observations;
they do not verify Galois descent on varieties or the analytic theorem.
The latter is the argument above, relative to the permitted theorem B
and ordinary stated background.

\texttt{build\_deliverable.py} runs these checks, preserves their stdout
and stderr, generates the standalone \texttt{P16.4-paper.tex} from this
argument with the existing pandoc, and runs the existing pdfLaTeX twice
with shell escape disabled. It records command arguments, return codes,
logs, environment details, source/input hashes, and scope status. Run
from this workspace:

\begin{Shaded}
\begin{Highlighting}[]
\ExtensionTok{python3}\NormalTok{ finite\_descent\_checks.py}
\ExtensionTok{python3}\NormalTok{ build\_deliverable.py}
\end{Highlighting}
\end{Shaded}

No installations, heavy compute, external workers, account changes,
editor operations, private streams, raw SQLite, or other campaigns were
used. No supervised heavy-compute facility was exposed, and none was
needed for this proof and these small exact checks. The requested
external editor path was not accessed or changed. Execution success and
compilation are observations, not mathematical certification. This is an
unverified draft with an unconditional boundary conclusion relative to
the supplied complete B grant and explicit gaps for the entire original
problem.

\subsection{References}\label{references}

Complete B: the entire accepted mathematical body supplied in the
assignment, extracted as \texttt{grant-B.tex} with surrounding blank
whitespace normalized; the untouched \texttt{input-original.json}
preserves the original bytes. Only its actual ample subcase over
\(\bar k\) is used.

S. Boucksom and D. Eriksson, \emph{Spaces of norms, determinant of
cohomology and Fekete points in non-Archimedean geometry},
\href{https://arxiv.org/pdf/1805.01016v3}{arXiv:1805.01016v3}.

L. Pille--Schneider, \emph{Global pluripotential theory on hybrid
spaces}, \href{https://arxiv.org/pdf/2209.04879v2}{arXiv:2209.04879v2},
diagnostic scope check.

S. Boucksom and M. Jonsson, \emph{Measures of finite energy in
pluripotential theory: a synthetic approach},
\href{https://arxiv.org/pdf/2307.01697v1}{arXiv:2307.01697v1},
diagnostic only.

A. Chambert-Loir and A. Ducros, \emph{Formes différentielles réelles et
courants sur les espaces de Berkovich},
\href{https://arxiv.org/pdf/1204.6277v2}{arXiv:1204.6277v2}, documentary
and continuity-domain check only.

\end{document}
